%!TEX root = main.tex

\chapter{Limites e Continuidade}

\section{Topologia do Espaço Euclidiano}

\begin{definition}
    Sejam $a \in \R^n$ e $r \in \R_{>0}$.
        \begin{enumerate}[leftmargin=*, align=left, label=\textbf{(\alph*)}]
            \item A \emph{bola aberta de centro $a$ e raio $r$} é definida como
                \[
                    B_r(a) := \{x \in \R^n : \|x-a\| < r\}.
                \]
            \item A \emph{bola aberta furada de centro $a$ e raio $r$} é definida como
                \[
                    B^*_r(a) := \{x \in \R^n : 0 < \|x-a\| < r\}.
                \]
        \end{enumerate}
\end{definition}

\begin{corollary}
    Sejam $a \in \R^n$ e $r \in \R_{>0}$. Então $B^*_r(a) = B_r(a) \setminus \{a\}$.
\end{corollary}

\begin{proof}
\end{proof}

\begin{definition}
    Seja $A \subseteq \R^n$.
        \begin{enumerate}[leftmargin=*, align=left, label=\textbf{(\alph*)}]
            \item Dizemos que $a \in A$ é um \emph{ponto interior} de $A$ se existe $r \in \R_{>0}$ tal que $B_r(a) \subseteq A$.
            \item O \emph{interior} de $A$ é definido como
                \[
                    \op{int}{A} := \{ x \in A : \exists r (r \in \R_{>0} \land B_r(x) \subseteq A ) \}.
                \]
        \end{enumerate}
\end{definition}

\begin{definition}
    Dizemos que $A \subseteq \R^n$ é \emph{aberto} se $\op{int}{A} = A$.
\end{definition}

\begin{corollary}
    Toda bola aberta é aberta.
\end{corollary}

\begin{proof}
    \blackproof
\end{proof}

\begin{theorem}
    \leavevmode
        \begin{enumerate}[leftmargin=*, align=left, label=\textbf{(\alph*)}]
            \item $\emptyset$ e $\R^n$ são abertos.
            \item Toda união de abertos do $\R^n$ é aberta.
            \item Toda interseção finita de abertos do $\R^n$ é aberta.
        \end{enumerate}
\end{theorem}

\begin{proof}
\end{proof}

\begin{definition}
    Sejam $a \in \R^n$ e $r \in \R_{>0}$.
        \begin{enumerate}[leftmargin=*, align=left, label=\textbf{(\alph*)}]
            \item A \emph{bola fechada de centro $a$ e raio $r$} é definida como
                \[
                    B_r[a] := \{x \in \R^n : \|x-a\| \leq r\}.
                \]
            \item A \emph{bola fechada furada de centro $a$ e raio $r$} é definida como
                \[
                    B^*_r[a] := \{x \in \R^n : 0 < \|x-a\| \leq r\}.
                \]
        \end{enumerate}
\end{definition}

\begin{definition}
    Dizemos que $A \subseteq \R^n$ é \emph{fechado} se $A^C$ é aberto.
\end{definition}

\begin{corollary}
    Toda bola fechada é fechada.
\end{corollary}

\begin{proof}
    \blackproof
\end{proof}

\begin{theorem}
    \leavevmode
        \begin{enumerate}[leftmargin=*, align=left, label=\textbf{(\alph*)}]
            \item $\emptyset$ e $\R^n$ são fechados.
            \item Toda união finita de fechados do $\R^n$ é fechada.
            \item Toda interseção de fechados do $\R^n$ é fechada.
        \end{enumerate}
\end{theorem}

\begin{proof}
    \blackproof
\end{proof}

\begin{definition}
    Seja $A \subseteq \R^n$.
        \begin{enumerate}[leftmargin=*, align=left, label=\textbf{(\alph*)}]
            \item Dizemos que $a \in \R^n$ é um \emph{ponto exterior} de $A$ se $a \in \op{int}{A^C}$.
            \item O \emph{exterior} de $A$ é definido como $\op{ext}{A} := \op{int}{A^C}$.
        \end{enumerate}
\end{definition}

\begin{definition}
    Seja $A \subseteq \R^n$.
        \begin{enumerate}[leftmargin=*, align=left, label=\textbf{(\alph*)}]
            \item Dizemos que $a \in \R^n$ é um \emph{ponto da fronteira} de $A$ se $a \notin \op{int}{A}$ e $a \notin \op{ext}{A}$.
            \item A \emph{fronteira} de $A$ é definida como
                \[
                    \partial A := \{ x \in \R^n : x \notin \op{int}{A} \land x \notin \op{ext}{A} \}.
                \]
        \end{enumerate}
\end{definition}

\begin{definition}
    Sejam $a \in \R^n$ e $r \in \R_{>0}$. A \emph{esfera de centro $a$ e raio $r$} é definida como
        \[
            S_r(a) := \{x \in \R^n : \|x-a\| = r\}.
        \]
\end{definition}

\begin{proposition}
    Sejam $a \in \R^n$ e $r \in \R_{>0}$.
        \begin{enumerate}[leftmargin=*, align=left, label=\textbf{(\alph*)}]
            \item $B_r[a] = B_r(a) \cup S_r(a)$.
            \item $\partial (B_r(a)) = \partial (B_r[a]) = S_r(a)$.
        \end{enumerate}
\end{proposition}

\begin{proof}
    \blackproof
\end{proof}

\begin{definition}
    Seja $A \subseteq \R^n$.
        \begin{enumerate}[leftmargin=*, align=left, label=\textbf{(\alph*)}]
            \item Dizemos que $a \in \R^n$ é um \emph{ponto de acumulação} de $A$ se
                \[
                    B^*_r(a) \cap A \neq \emptyset
                \]
            para todo $r \in \R_{>0}$.
            \item O conjunto de todos os pontos de acumulação de $A$ é denotado por $A'$, isto é,
                \[
                    A' := \{x \in \R^n : \forall r (r \in \R_{>0} \rightarrow B^*_r(x) \cap A \neq \emptyset) \}.
                \]
        \end{enumerate}
\end{definition}

\begin{definition}
    Dizemos que $a \in A$ é um \emph{ponto isolado} de $A \subseteq \R^n$ se $a \notin A'$.
\end{definition}

\begin{definition}
    Seja $A \subseteq \R^n$.
        \begin{enumerate}[leftmargin=*, align=left, label=\textbf{(\alph*)}]
            \item Dizemos que $a \in \R^n$ é um \emph{ponto de aderência} de $A$ se
                \[
                    B_r(a) \cap A \neq \emptyset
                \]
            para todo $r \in \R_{>0}$.
            \item O \emph{fecho} de $A$ é definido como
                \[
                    \bar{A} := \{x \in \R^n : \forall r (r \in \R_{>0} \rightarrow B_r(x) \cap A \neq \emptyset) \}.
                \]
        \end{enumerate}
\end{definition}

\begin{proposition}
    Seja $A \subseteq \R^n$. 
        \begin{enumerate}[leftmargin=*, align=left, label=\textbf{(\alph*)}]
            \item 
                \begin{enumerate}[label=\roman*.]
                    \item $\op{int}{A}$ é aberto.
                    \item $\bar{A}$ é fechado.
                    \item $\partial A$ é fechado.
                \end{enumerate}
            \item 
                \begin{enumerate}[label=\roman*.]
                    \item $\op{int}{A} \cap \partial A = \emptyset$.
                    \item $\bar{A} = A \cup A'$.
                    \item $\bar{A} = A \cup \partial A$.
                \end{enumerate}
            \item Se $a \in A'$, então $a \in \op{int}{A}$ ou $a \in \partial A$.
            \item $A$ é fechado se, e somente se, $\bar{A} = A$.
        \end{enumerate}
\end{proposition}

\begin{proof}
    \blackproof
\end{proof}

\begin{definition}
    Dizemos que $A \subseteq \R^n$ é \emph{limitado} se existe $r \in \R_{>0}$ tal que $A \subseteq B_r(0)$.
\end{definition}

\begin{definition}
    Dizemos que $A \subseteq \R^n$ é \emph{compacto} se é fechado e limitado.
\end{definition}

\section{Limites}

\begin{definition}
    Dizemos que $L \in \R^m$ é o \emph{limite} de $f : A \subseteq \R^n \to \R^m$ em $a \in A'$ se para todo $\epsilon \in \R_{>0}$ existe $\delta = \delta(\epsilon, a) \in \R_{>0}$ tal que, para todo $x \in A$, se $0 < \| x-a \| < \delta$, então $\| f(x)-L \| < \epsilon$. Isso é denotado por 
        \[
            \lim_{x \to a} f(x) = L.
        \]
\end{definition}

\begin{proposition}[Unicidade] \label{teo.calc2:unicidade_limite_bilateral}
    O limite de $f : A \subseteq \R^n \to \R^m$ em $a \in A'$, quando existe, é único.
\end{proposition}

\begin{proposition}
    Sejam $f : A \subseteq \R^n \to \R^m$, $a \in A'$ e $L \in \R^m$. Então
        \[
            \lim_{x \to a} f(x) = L \Leftrightarrow \lim_{\| x-a \| \to 0} \|f(x) - L \| = 0.
        \]
\end{proposition}

\begin{proof}
    \blackproof
\end{proof}

\section{Continuidade}

\begin{definition}[Continuidade] \label{def.calc2:continuidade}
    Dizemos que $f : A \subseteq \R^n \to \R^m$ é \emph{contínua em $a \in A$} se para todo $\epsilon \in \R_{>0}$ existe $\delta = \delta(\epsilon, a) \in \R_{>0}$ tal que, para todo $x \in A$, se $\| x-a \| < \delta$, então $\| f(x) - f(a) \| < \epsilon$.
\end{definition}

\begin{definition}
    Sejam $f : A \subseteq \R^n \to \R^m$ e $X \subseteq A$.
        \begin{enumerate}[leftmargin=*, align=left, label=\textbf{(\alph*)}]
            \item Dizemos que $f$ é \emph{contínua em $X$} se $f$ é contínua em todos os pontos de $X$.
            \item Dizemos que $f$ é \emph{contínua} se $f$ é contínua em $A$.
        \end{enumerate}
\end{definition}

\begin{theorem} \label{teo.calc2:continuidade_sse_limite}
    Sejam $f : A \subseteq \R^n \to \R^m$ e $a \in A \cap A'$. Então, $f$ é contínua em $a$ se, e somente se, $\ds \lim_{x \to a} f(x) = f(a)$.
\end{theorem}

\begin{proof}
    \blackproof
\end{proof}

\begin{proposition}
    Sejam $f : A \subseteq \R^n \to \R^m$ e $g : B \subseteq \R^m \to \R^k$ tais que $f(A) \subseteq B$. Se $f$ é contínua em $a \in A$ e $g$ é contínua em $f(a) \in B$, então $g \circ f : X \subseteq \R^n \to \R^k$ é contínua em $a$.
\end{proposition}

\begin{proof}
    \blackproof
\end{proof}